\section{Relation to Cognitive Theories and Testable Hypotheses}
\label{sec:cognitive_systems}

\subsection{Cognitive interpretation}
\label{subsec:cognition_forces}
For a cognitive implementation, the state might include a belief representation, a policy, retained knowledge, and resource estimates. Epistemic pull then concerns informative sampling; selective protection concerns retention and integrity; computational drag concerns the cost of inference; action efficiency concerns task performance and timeliness; and sustainability concerns future capacity to continue those activities. These are modelling assignments whose separability must be tested.

A cognitive Variational Equilibrium uses the same stationary definition as Eq.~\eqref{eq:veq}. It is not defined as a point where one objective cannot improve without harming another: that would introduce a different, Pareto-based notion. Continuing cognition is more naturally described through trajectories in an admissible operating set, with equilibria relevant only to particular state summaries or fixed conditions.

\subsection{Predictive processing and active inference}
\label{subsec:pp_ai}
Predictive processing emphasizes hierarchical predictions and the updating of beliefs through prediction errors \citep{clark2013whatever}. Active inference provides a generative-model account of perception and action, including epistemic and pragmatic contributions to policy evaluation \citep{friston2010fep,parr2018generalised,parr2022active}. Combining these contributions in one functional does not make them indistinguishable or imply that their behavioural balance must be fixed. That balance depends on beliefs, preferences, policies, and the model's treatment of the environment and internal states.

\uvif{} makes a different organizational choice by assigning named update components to regulatory demands. This may facilitate intervention and reporting, but it is not yet a derivation of a distinct class of behaviour. In particular, a resource-induced change in exploration is not by itself evidence against active inference: an active-inference model that represents resource state may also predict it. A valid comparison must identify which variables and mechanisms each fitted model includes.

Likewise, a common active-inference formulation may omit a detailed account of inference expenditure without that omission being an impossibility of the framework. The pertinent comparison is between specified models with and without explicit computational and reserve accounting. Claims about entire theoretical families would exceed what the present framework establishes.

\subsection{Resource-rational analysis}
\label{subsec:resource_rational}
Resource-rational analysis explains cognitive strategies in terms of expected benefits relative to computational costs \citep{lieder2019resourcerational,griffiths2015rational}. Human planning studies provide examples of how such models can be tested \citep{callaway2022rational}. This literature directly anticipates the interaction between action efficiency and computational drag.

The proposed sustainability component asks whether the resource state carried between decisions changes the preferred computation policy. A per-task model that omits those dynamics may fail, but an extended resource-rational model can represent them. Irreversible outcomes can similarly be represented through absorbing states, constraints, or preferences. UVIF's possible contribution is a convenient decomposition and independently measurable mechanisms, not the exclusive ability to represent these considerations.

\subsection{Allostasis and cognitive architectures}
\label{subsec:allostasis}
Allostasis concerns predictive physiological regulation \citep{schulkin2019allostasis}. Evidence for a large-scale brain system supporting allostasis and interoception motivates examining cognition and bodily regulation together \citep{kleckner2017evidence}. It does not establish a one-to-one mapping between a neural network and the five proposed forces. Such a mapping would require independent measurements and interventions.

\label{subsec:architectures}
Integrated cognitive architectures describe mechanisms for memory, learning, control, and action \citep{laird2017standard,kotseruba201840}. \uvif{} is a proposed regulatory description that could be implemented within such architectures. It does not specify a complete cognitive architecture. A concrete implementation must explain how component signals are computed, how competing updates are resolved, and how the computation used by the regulatory layer itself is charged to the resource budget.

\subsection{Comparative hypotheses}
\label{subsec:cog_predictions}
The following hypotheses concern specified implementations rather than unique predictions of UVIF as a whole. Their value lies in enabling controlled comparisons and revealing when the additional decomposition is unnecessary.

\begin{enumerate}[label=(H\arabic*),leftmargin=2.4em]
\item \textbf{Resource-dependent allocation.} Manipulating available computation, time, or reserve state while preserving external task information and rewards should alter information acquisition when the specified model couples allocation to those resources. Compare a resource-responsive UVIF model with both fixed-weight and resource-responsive alternatives. Measure sampling decisions, response time, task performance, and resource expenditure. A failure to improve held-out prediction over equally flexible alternatives would weaken the case for the proposed decomposition.

\item \textbf{Effects of resource carry-over.} When expenditure reduces capacity in later episodes, a model that tracks reserves should predict different allocation from an otherwise comparable model that resets resources after each episode. Match the observed task and state information across agents, and compare with a long-horizon resource-rational or constrained-control model. The prediction concerns explicit carry-over dynamics, not an assumed distinction between a sustainability force and all possible long-horizon utilities.

\item \textbf{Protection under irreversible outcomes.} Compare a hard admissibility mechanism with calibrated soft penalties under the same transition model, including irreversible states. Measure violation frequency, retained task performance, and unnecessary refusal. Hard and soft formulations may coincide over a bounded range of payoffs; the experiment must report that range. No behavioural observation here establishes impossibility of representation by every utility model.

\item \textbf{Component ablation.} Remove one implemented component while preserving observations and evaluating both fixed-parameter and retuned controls. Predefine the expected consequence: reduced learning under epistemic ablation, greater damage under protection ablation, or reserve depletion under sustainability ablation. Effects need not be monotonic or unique, because other components may compensate. An ablation is informative only when it tests an independently meaningful mechanism rather than a failure built into the outcome definition.
\end{enumerate}

Section~\ref{sec:future} specifies common evaluation requirements. This paper does not report tests of H1--H4 or claim that they distinguish UVIF from every member of the comparator families. The companion repository provides an executed notebook and saved outputs for illustrative policy-grid demonstrations \citep{Hamam2026GitHub}. Its implementation maximizes a scalar score formed from the five named contributions, rather than integrating Eq.~\eqref{eq:projected}, and the notebook records two negative results that bear directly on the hypotheses above. First, a fixed unit-weight scalarization over the same state-dependent component functions reproduces every policy that implementation selects, across all declared environments; the implementation is therefore a scalar optimization and cannot evidence a departure from one. Second, a finite monotone hinge penalty on exposure above the admissibility threshold reproduces the constrained policies exactly over the same declared set, which refutes any claim that constrained choice under irreversibility is irreducible to a finite monotone penalty. Both results are exhaustive over the declared environment set and policy grid. They are reported as demonstrations of what that implementation cannot establish, and they motivate the evaluation requirements of Section~\ref{sec:future}; neither the notebook nor this paper validates H1--H4 or the vector-field model.
